\scanpage{74}
% manuscript line 1
coïncident dans les $T_i\cap T_j$, donc définissent $T\unsure{\text{---}|\longrightarrow}X$ et ce dernier satisfait aux conditions voulues. D'ailleurs, \deleted{\unsure{toute|t.t|tout}} \unsure{une} solution ne peut, \unsure{à|a} priori, s'obtenir par une construction comme celle qui précède. Introduisant alors les
% manuscript line 7
pour $u,u'$ \unsure{liés|choisis}\,: $(T_i)$, $(T'_j)$, le \unsure{recouvrt|recouvrement} commun $(T_i\cap T'_j)_{i,j\in I\times I'}$), on trouve l'unicité.

% manuscript line 10
\uline{Lemme 3} La condition $\beta$ est automatiquement vérifiée \inserted{(pour tt rec. $\unsure{\mathfrak X|\mathscr X}=\bigcup\mathfrak X_i$ de $\mathfrak X$ par ouverts)} si on a un isomorphisme de foncteurs
% manuscript line 14
\[
|T|\overset{\sim}{\unsure{\longrightarrow|\text{---}}}|\underline W(T)|
\]
% manuscript line 15
en particulier dans les cas $\mathrm{II}'$, \deleted{$\unsure{\mathrm{II}|\mathrm{I}}$}, en particulier $\mathrm{III}$.

% manuscript line 16
C'est trivial.

% manuscript line 17
\uline{Lemme 4} La condition $\alpha$ est automatiquement vérifiée dans les cas $\mathrm{I}',\mathrm{II}^{*}$.

% manuscript line 19
C'est trivial pour $\mathrm{II}^{*}$. Pour $\mathrm{I}'$, petit raisonnement, utilisant que un morphisme fini est \uline{fermé}. (\unsure{Les hyp. inutiles|Sans hyp. noeth.})
% archive sheet number: 73
